% Arquivo: latex/politica-regional.tex
% Seção 2: Política Regional no Brasil
% Artigo: Avaliação dos Instrumentos da PNDR
% Gerado em: 2026-03-02

\section{Política Regional no Brasil}
\label{sec:politica_regional}

A desigualdade regional constitui um dos traços estruturais mais persistentes da economia brasileira. Apesar dos avanços sociais e econômicos das últimas décadas, a maioria dos municípios das áreas de atuação da SUDAM, SUDENE e SUDECO ainda apresenta PIB \textit{per capita} inferior à média nacional, com lenta convergência ao longo do período 2002--2021 \cite{Portugal2024}. Essa persistência motiva a existência de políticas específicas de desenvolvimento regional, cujas origens, institucionalização e instrumentos são descritos nesta seção.

\subsection{Origens e evolução histórica da política regional}
\label{subsec:origens_evolucao}

A região Nordeste foi a primeira macrorregião reconhecida como ``região-problema'', termo cunhado pelo Grupo de Trabalho para o Desenvolvimento do Nordeste (GTDN) para referir-se ao seu nível de desigualdade e atraso econômico em relação ao Centro-Sul \cite{Portugal2024}. Para \citeonline{Moreira2014}, a problemática da desigualdade regional emerge da divisão territorial do trabalho resultante do processo de industrialização nacional entre 1880 e 1930, intensificando-se após 1950. O contraste entre as estruturas produtivas --- o Nordeste açucareiro/pecuário-algodoeiro e o Centro-Sul cafeeiro/industrial --- marcou profundamente as relações de produção entre as principais regiões do país.

A linha teórico-estruturalista de \citeonline{Prebisch1949}, que tem em \citeonline{Furtado2009} seu maior expoente no país, explica as desigualdades regionais pela reprodução, em âmbito nacional, da relação de dependência centro-periferia característica da economia mundial: metrópoles industrializadas produzindo bens de alto valor agregado versus periferias produtoras de matérias-primas, condenando estas últimas ao subdesenvolvimento.

Os primeiros esforços sistemáticos de política de desenvolvimento regional no Brasil remontam às décadas de 1950 e 1960, com a criação da Superintendência de Valorização Econômica da Amazônia (SPVEA), pela Lei nº 1.806/1953, e da SUDENE, pela Lei nº 3.692/1959 \cite{Silva2014, PortugalSilva2020}. A SUDENE, instituída sob a visão teórica de Celso Furtado, tornou-se central no esforço estratégico do governo federal para superar o atraso relativo da região, partindo do diagnóstico de que era preciso industrializá-la por meio da participação ativa do Estado. Durante o governo de Juscelino Kubitschek (1956--1961), um arranjo federativo atribuiu certa autonomia às superintendências regionais, possibilitando diálogo com os governos estaduais e fortalecendo o federalismo cooperativo \cite{Bercovici2003}.

O principal instrumento de atuação, entretanto, não foi a ação direta do Estado via investimento público, mas a concessão de subsídios fiscais e financeiros a empresas que realizassem investimentos produtivos nas regiões prioritárias, conforme as Leis nº 4.069-B/1962 (isenção de IR na Amazônia) e nº 4.239/1963 (isenção e dedução de IR no Nordeste, com instituição do FIDENE).

Durante o Regime Militar (1964-1985), o modelo existente no Nordeste foi reproduzido em outras macrorregiões: a SPVEA foi substituída pela SUDAM (Lei nº 5.173/1966), criou-se a SUFRAMA para a Zona Franca de Manaus, e a antiga Fundação Brasil-Central deu lugar à SUDECO (Lei nº 5.365/1967). Não obstante a expansão institucional, o período foi caracterizado pela centralização do planejamento governamental, retirando autonomia política e administrativa dos órgãos regionais. \citeonline{Silva2014} aponta incremento das políticas de ordenamento territorial e declínio das políticas regionais baseadas em diagnósticos de especificidades locais, com o planejamento regional tornando-se apenas ``capítulos'' de uma estratégia nacional após o lançamento dos Planos Nacionais de Desenvolvimento. Em 1974, foram instituídos os Fundos de Investimento da Amazônia (FINAM) e do Nordeste (FINOR), via Decreto-Lei nº 1.376/1974, instrumentos que permitiam aplicação de parcela do IR em debêntures conversíveis emitidas por empreendimentos estratégicos.

A partir do final dos anos 1970, o esgotamento da capacidade de financiamento do Estado brasileiro aprofundou o esvaziamento das instituições regionais e sua subordinação ao planejamento econômico nacional \cite{Silva2014}. As décadas de 1980 e 1990 foram marcadas por reformulação profunda do papel do Estado na economia, em meio a sucessivos planos de estabilização fiscal e monetária, com o Estado desenvolvimentista perdendo espaço para o Estado facilitador, regulador e fiscalizador das atividades econômicas privadas. Nesse contexto, o FINAM e o FINOR foram revogados em 2001, em meio a denúncias de improbidade \cite{Portugal2024}. A retomada da questão regional, contudo, já se delineava desde a promulgação da Constituição Federal de 1988.

\subsection{Política Nacional de Desenvolvimento Regional (PNDR)}
\label{subsec:pndr}

A Constituição Federal de 1988 estabeleceu as bases para a retomada da questão regional ao inserir a redução das desigualdades regionais como princípio fundamental (Art. 3º, inciso III) e ao prever a participação colaborativa da União nos problemas regionais (Art. 43). O novo texto constitucional modificou substancialmente a estrutura federativa, aumentando a participação de estados e municípios na arrecadação federal. De acordo com o Art. 159, do montante arrecadado com impostos sobre a renda e produtos industrializados, 21,5\% seriam destinados ao Fundo de Participação dos Estados, 22,5\% ao Fundo de Participação dos Municípios e 3\% a programas de financiamento produtivo de caráter regional. Para \citeonline{PortugalSilva2020}, essa estrutura imprensou as entidades regionais entre a União e os Estados no novo arranjo federativo, dado que a maior parte dos recursos segue diretamente para estados e municípios.

Seguindo o disposto na CF/1988, foram criados os Fundos Constitucionais de Financiamento para as regiões Norte (FNO), Nordeste (FNE) e Centro-Oeste (FCO), regulamentados pela Lei nº 7.827/1989, com o objetivo de direcionar crédito aos setores produtivos, preferencialmente para mini e pequenos produtores. A repartição dos 3\% da arrecadação é feita na forma de 1,8\% para o FNE (dos quais 0,9\% deve ser obrigatoriamente aplicado à região do semiárido), 0,6\% para o FNO e 0,6\% para o FCO.

No contexto de reforma da estrutura estatal, em 2001 as superintendências SUDAM e SUDENE foram extintas por meio da MP nº 2.145/2001, sendo substituídas pelas Agências de Desenvolvimento da Amazônia e do Nordeste (ADA e ADENE). Para \citeonline{PortugalSilva2020}, os novos órgãos passaram a atuar essencialmente como executores de planos de desenvolvimento, e não como formuladores. A mesma MP instituiu os Fundos de Desenvolvimento da Amazônia (FDA) e do Nordeste (FDNE), fundos contábeis cujo objetivo é assegurar recursos para o financiamento de investimentos privados de grande porte e considerados estratégicos ao desenvolvimento regional.

Em 2003, com o início do governo Lula, houve esforço para restabelecer o protagonismo das políticas regionais no país, sob a liderança do Ministério da Integração Nacional. Os debates resultaram, em 2007, na promulgação do Decreto nº 6.047/2007, que instituiu a PNDR, com o objetivo expresso de reduzir as desigualdades de qualidade de vida entre as regiões brasileiras e promover a equidade no acesso a oportunidades de desenvolvimento. A PNDR surge do diagnóstico de que ``a desigualdade regional é resultado da dinâmica assimétrica do crescimento capitalista, que se concentra em alguns espaços enquanto condena outros à estagnação'' \cite{BrasilPndr2006}.

A PNDR tem três premissas fundamentais: (i) é uma política nacional, no sentido de que compete ao Governo Federal emprestar coerência e efetividade aos esforços de desenvolvimento regional; (ii) é uma política do governo federal, devendo envolver, em graus variados, todas as áreas de atuação do governo; (iii) exige uma abordagem das desigualdades regionais em múltiplas escalas, reconhecendo diferentes dinâmicas sub-regionais \cite{BrasilPndr2006}.

A fim de orientar as ações da política, foi estabelecida uma tipologia técnica de áreas prioritárias baseada em regiões geográficas imediatas. O primeiro quadro de tipologias, definido pelo Decreto nº 6.047/2007, baseou-se na combinação do rendimento médio mensal por habitante e da taxa geométrica de variação do PIB \textit{per capita}, considerando prioritárias as sub-regiões classificadas como de baixa renda, estagnadas ou dinâmicas. Ainda em 2007, as superintendências regionais SUDAM e SUDENE foram restabelecidas por meio das Leis Complementares nº 124/2007 e nº 125/2007, retomando o lugar que haviam perdido para as agências, e em 2009 a SUDECO foi instituída pela Lei Complementar nº 129/2009. Essas instituições ficaram responsáveis pela gestão e articulação da PNDR, visando à convergência das políticas públicas dos diferentes níveis de governo. Como forma de financiamento, a PNDR conta com recursos do orçamento geral da União e dos entes federativos, os Fundos Constitucionais de Financiamento, Fundos de Desenvolvimento Regional, Incentivos Fiscais e o Fundo Nacional de Desenvolvimento Regional (FNDR), este último regulamentado apenas em 2023 (Emenda Constitucional nº 132/2023).

Não obstante os avanços institucionais, \citeonline{PortugalSilva2020} avaliam que a política falhou em seus objetivos, devido à: (i) inexistência de marco normativo com maiores vínculos jurídicos (como uma lei, em lugar de decreto); (ii) baixa colaboração e coordenação entre entes federativos e ministérios; (iii) a questão regional não entrou na agenda de prioridades do governo, e os planos regionais não foram instituídos. \citeonline{Coelho2015} reforça que as instituições regionais não contaram com instrumentos efetivos para enfrentar as desigualdades regionais e lembra que a operacionalização da política foi comprometida pela não aprovação do FNDR.

Em 2013, uma nova versão da PNDR foi elaborada no âmbito da I Conferência Nacional do Desenvolvimento Regional, sendo formalizada em 2019 pelo Decreto nº 9.810/2019 (PNDR II). A PNDR II estabelece quatro objetivos estratégicos: (i) promover a convergência dos níveis de desenvolvimento econômico e qualidade de vida entre e intrarregiões; (ii) consolidar uma rede policêntrica de cidades; (iii) estimular a competitividade em regiões de declínio econômico e populacional; (iv) fomentar a agregação de valor em regiões de alta especialização em \textit{commodities} agrominerais. Como estratégia de ação, a PNDR II promove a estruturação do Sistema de Governança do Desenvolvimento Regional e o estabelecimento de seis eixos setoriais de atuação: desenvolvimento produtivo; ciência, tecnologia e inovação; educação e qualificação profissional; infraestrutura econômica e urbana; desenvolvimento social; e fortalecimento das capacidades governativas dos entes federativos.

Em 2018, foi realizada a atualização da tipologia municipal (Portaria nº 34/2018-MIN), com revisão dos indicadores, substituição do rendimento domiciliar médio pela paridade do poder de compra e utilização de médias móveis trienais para a variável dinâmica do PIB municipal. A nova tipologia classifica os municípios em seis grupos, conforme o cruzamento da variável rendimento/habitante (alto, médio e baixo) e a variação do PIB/habitante (alta, média e baixa).

\subsection{Fundos Constitucionais de Financiamento}
\label{subsec:fcs}

O objetivo dos Fundos Constitucionais FNO, FNE e FCO é contribuir para o desenvolvimento econômico e social das regiões de competência mediante a execução de programas de crédito subsidiado aos setores produtivos locais. A Lei nº 7.827/1989, que regulamenta esses fundos, estipula tratamento preferencial a pequenos e miniprodutores rurais, pequenas e microempresas; prestação de assistência técnica em caso de setores carentes; e apoio à formação de novos centros, atividades e polos dinâmicos. Um aspecto importante é que a previsão constitucional dos FCs confere a eles um fluxo seguro de recursos financeiros, sem contingenciamentos e livre de fatores conjunturais políticos e econômicos, diferenciando-os dos demais instrumentos de política regional. 

Os FCs mobilizam o maior volume de recursos dentre os instrumentos da PNDR. Em 2020, foram aplicados R\$ 25,8 bilhões pelo FNE, R\$ 10,5 bilhões pelo FNO e R\$ 7,5 bilhões pelo FCO. Conforme a Tabela~\ref{tab:resumo_fc}, os valores aplicados mais que duplicaram entre o período 2002-2008 e 2009-2015, sendo mais expressivo nos municípios de Alta-Renda atendidos pelo FNO (192\%) e nos municípios de tipologia Dinâmica do FNE (151\%). Entre os períodos 2009-2015 e 2016-2021 o desempenho é mais moderado, o contexto de menor crescimento nacional significou redução das receitas dos fundos. Esse refletiu em menos recursos aos municípios de Alta-Renda, em especial no FNO (-32,8\%). Para \citeonline{PortugalSilva2020}, os FCs são um exemplo em que ``o instrumento se tornou maior do que a política'', uma vez que a PNDR, instituída por decretos presidenciais, possui menor segurança financeira e jurídica que os Fundos Constitucionais, que tramitaram no Congresso Nacional e foram sancionados na forma de lei.

\begin{table}[h!]
	\centering
	\caption{Aplicações dos Fundos Constitucionais (R\$ de 2020)}
	\label{tab:resumo_fc}
	\footnotesize
	\renewcommand{\arraystretch}{1.2}
	\begin{tabular}{llcccccc}
		\toprule
		\multirow{2}{*}{Fundo} & \multirow{2}{*}{Tipologia} & \multicolumn{3}{c}{Valor total (R\$ Milhões)} & \multicolumn{3}{c}{Média \textit{per capita}} \\
		\cmidrule(lr){3-5} \cmidrule(lr){6-8}
		& & 2002-2008 & 2009-2015 & 2016-2021 & 2002-2008 & 2009-2015 & 2016-2021 \\
		\midrule
		\multirow{4}{*}{FNE} & Alta Renda & 12.441 & 26.441 & 66.420 & 23,20 & 9,54 & 7,93 \\
		& Baixa Renda & 8.329 & 26.151 & 52.095 & 16,11 & 13,01 & 12,04 \\
		& Dinâmica & 14.135 & 40.330 & 90.008 & 30,80 & 34,71 & 30,04 \\
		& Estagnada & 17.895 & 54.073 & 114.617 & 18,69 & 15,03 & 14,22 \\
		\midrule
		\multirow{3}{*}{FCO} & Alta Renda & 12.788 & 19.981 & 58.152 & 22,26 & 13,32 & 17,28 \\
		& Dinâmica & 4.567 & 11.163 & 26.535 & 26,01 & 32,19 & 32,73 \\
		& Estagnada & 8.556 & 18.077 & 45.050 & 20,78 & 20,66 & 23,99 \\
		\midrule
		\multirow{4}{*}{FNO} & Alta Renda & 4.941 & 9.689 & 29.054 & 49,26 & 43,99 & 34,88 \\
		& Baixa Renda & 1.750 & 3.431 & 8.332 & 62,31 & 70,74 & 55,80 \\
		& Dinâmica & 3.536 & 10.643 & 20.267 & 48,14 & 54,04 & 50,69 \\
		& Estagnada & 7.311 & 17.128 & 38.372 & 83,28 & 79,60 & 75,12 \\
		\bottomrule
		\multicolumn{8}{l}{\footnotesize{Fonte: Elaborada pelo autor a partir de dados de BB, BASA e BNB.}} \\
	\end{tabular}
\end{table}

A participação dos recursos no PIB local, apresentada na Tabela~\ref{tab:resumo_fc}, revela padrão relevante de concentração. No FNE, municípios de tipologia dinâmica apresentam a maior participação média no PIB anual (6,14\% no subperíodo 2016--2021), seguidos pelos estagnados (3,55\%) e de baixa renda (3,16\%), enquanto municípios de alta renda registram participação de apenas 1,12\%. A trajetória temporal indica crescimento expressivo da participação nos municípios dinâmicos --- de 2,38\% em 2002--2008 para 6,14\% em 2016--2021 --- e de baixa renda (de 1,77\% para 3,16\%), sugerindo intensificação recente das aplicações nesses grupos.

O FCO apresenta dinâmica semelhante, com maior participação em municípios dinâmicos (4,68\%) e estagnados (4,22\%), e redução relativa nos de alta renda (de 2,22\% para 2,49\%). No FNO, a participação é mais elevada em municípios estagnados (5,43\%) e dinâmicos (3,68\%), refletindo as características da ocupação territorial na região Norte. De modo geral, a distribuição relativa dos FCs revela concentração de recursos em municípios de tipologia dinâmica e estagnada, com participação proporcionalmente menor em municípios de alta renda. Considerando que a PNDR tem como objetivo prioritário a redução das desigualdades regionais, essa distribuição sugere alinhamento parcial dos FCs com as prioridades da política, embora a concentração em municípios dinâmicos --- e não nos de baixa renda --- levante questões sobre a efetiva focalização dos recursos.

\subsection{Fundos de Desenvolvimento}
\label{subsec:fds}

FDA, FDNE e FDCO são fundos de natureza contábil que visam financiar projetos de investimento em infraestrutura e capital fixo de alto potencial gerador de novos negócios e atividades produtivas. Diferentemente dos FCs, os Fundos de Desenvolvimento focam exclusivamente em projetos de grande porte (mínimo de R\$ 15 milhões) e não operam diretamente com o beneficiário final. Compete às superintendências SUDAM, SUDENE e SUDECO administrar os respectivos fundos, analisar a adequação dos pedidos de apoio financeiro, celebrar contratos com os agentes financeiros operadores e auditar o uso dos recursos. As aplicações devem seguir as diretrizes e prioridades estabelecidas nos planos regionais e sub-regionais de desenvolvimento, pela PNDR e pelo Ministério do Desenvolvimento Regional.

A operacionalização dos Fundos de Desenvolvimento utiliza um modelo de intermediação bancária obrigatória, contrastando com o acesso direto dos Fundos Constitucionais. Uma instituição bancária, denominada agente operador, faz a intermediação entre o empreendedor e o fundo, sendo responsável por fiscalizar e atestar as informações prestadas pelo proponente, avaliar a viabilidade econômico-financeira do projeto, negociar os aspectos de contratação e assumir o risco de crédito em cada operação. Esse modelo implica processo mais complexo e criterioso de aprovação, adequado para investimentos estruturantes de maior valor e impacto. As condições de financiamento variam conforme modalidade, setor e região de instalação, com períodos de até 20 anos e participação mínima própria de 20\%.

Segundo dados das superintendências, o FDNE executou R\$ 11,6 bilhões em financiamentos no período de 2008 a 2021, com cerca de 70\% alocados em projetos de infraestrutura, 29\% na indústria de transformação e o restante em indústria extrativa. Municípios classificados como estagnados respondem por 42\% do valor total (R\$ 4,9 bilhões), concentrando quase a totalidade dos empreendimentos em indústria de transformação. Esses municípios são os maiores beneficiários do FDNE, seja em termos absolutos ou relativamente ao tamanho da economia local. Dentre os projetos financiados está o polo automotivo em Goiana/PE, a Ferrovia Transnordestina e um total de 21 parques de geração de energia eólica distribuídos nos estados de Ceará, Rio Grande do Norte e Bahia. Alguns estudos já têm apontado efeitos econômicos significativos desses investimentos: \citeonline{Costaetal2024} verificaram efeito causal do Polo Automotivo de Goiana/PE sobre o PIB \textit{per capita} do município, e \citeonline{Irffietal2025} observaram impactos positivos da instalação de parques eólicos sobre o desenvolvimento local em municípios do Nordeste brasileiro.

Os financiamentos do FDA envolvem cerca de R\$ 5,5 bilhões, de 2007 a 2015, sendo 95\% aplicados em projetos de infraestrutura. Municípios de Alta Renda têm a maior participação (36\%), seguidos pela tipologia dinâmica (27\%), mas se comparado à magnitude da economia local, os recursos são mais significativos em municípios de baixa renda, dinâmicos e estagnados. O FDCO tem liberações de recursos apenas a partir de 2014, totalizando R\$ 1,8 bilhão, sendo direcionados principalmente à indústria de transformação (85\%) e municípios de alta renda (61\%). A Figura~\ref{fig:fd_setor} sintetiza o valor total liberado por cada fundo segundo o setor econômico dos projetos financiados.

% Gerado por: scripts/generate_policy_figures.py
\begin{figure}[htbp]
	\centering
	\caption{Valor liberado pelos Fundos de Desenvolvimento, por tipologia municipal e setor (2007-2021)}
	\label{fig:fd_setor}
	\includegraphics[width=1\textwidth]{fd_fundo_setor.png}
	\fonte{Elaborada pelos autores.}
\end{figure}

\subsection{Incentivos Fiscais}
\label{subsec:ifs}

Os Incentivos Fiscais são os instrumentos de desenvolvimento regional mais longevos, tendo sido promulgados pela Lei nº 4.069/1962 no âmbito da Amazônia Legal e pela Lei nº 4.239/1963 na região Nordeste, estabelecendo redução ou isenção de IR às empresas instaladas nessas regiões. Diferentemente dos Fundos Constitucionais e dos Fundos de Desenvolvimento, que envolvem financiamento direto, os Incentivos Fiscais operam por meio de renúncia tributária, ou seja, redução da carga de impostos devidos pelas empresas beneficiadas. Em 1997, a Lei nº 9.532/1997 estipulou a diminuição gradativa do percentual de redução fiscal até sua completa extinção em 2013. Porém, a Medida Provisória nº 2.199-14/2001 revogou a regulamentação anterior e firmou uma redução fixa de 75\% do IRPJ para empresas beneficiadas, inicialmente válida até 2018, posteriormente postergada até 2023 e, recentemente, estendida até 2028 pela Lei nº 14.753/2023.

Segundo \citeonline{ShirasuCorreaIrffi2023} e \citeonline{Portugal2024}, a hipótese subjacente ao programa é que a redução do custo tributário gera incentivos para empresas investirem e se instalarem em regiões de menor desenvolvimento, atuando sobre o custo de oportunidade do capital e tornando mais atrativo investir nas regiões beneficiadas. São elegíveis aos benefícios os projetos de instalação, ampliação, modernização ou diversificação de linhas de produção de empresas enquadradas em setores classificados como prioritários para o desenvolvimento regional nas áreas de atuação da SUDENE e SUDAM. Porém, \citeonline{Portugal2024} aponta que há uma desvinculação dos Incentivos Fiscais em relação aos objetivos do desenvolvimento regional e defende uma redefinição completa do programa, argumentando que a definição dos setores prioritários é prerrogativa do presidente da República (Decreto nº 4.212/2002), apesar da existência dos planos regionais de desenvolvimento elaborados em discussão com a sociedade.

Dados de renúncia fiscal da Receita Federal mostram que em 2021 as renúncias de IRPJ decorrentes dos incentivos da SUDAM somaram cerca de R\$ 12,6 bilhões, sendo R\$ 9,7 bilhões ligados à indústria e R\$ 2,8 bilhões à agricultura. No âmbito da SUDENE, o montante alcançou R\$ 19,3 bilhões, sendo R\$ 14,9 bilhões destinados à indústria e R\$ 4,4 bilhões à agricultura. Em um período mais longo, de 2015 a 2024, os dois programas de incentivo fiscal respondem por R\$ 126 bilhões em gastos tributários. A Figura~\ref{fig:incentivos} apresenta a distribuição dos benefícios por superintendência, setor e tipologia municipal.

% Gerado por: scripts/generate_policy_figures.py
\begin{figure}[htbp]
	\centering
	\caption{Distribuição dos benefícios de Incentivo Fiscal, por superintendência, setor e tipologia municipal (2010-2021)}
	\label{fig:incentivos}
	\includegraphics[width=1\textwidth]{icf_superint_setor.png}
	\fonte{Elaborada pelos autores.}
\end{figure}

No período de 2010 a 2021, SUDAM e SUDENE apoiaram 1.851 e 2.590 empresas com Incentivos Fiscais, respectivamente. Os projetos da SUDAM apresentam alta concentração em municípios de alta renda, com 1.260 empresas, contra apenas 34 em municípios de baixa renda. Na SUDENE, os incentivos se distribuem de forma mais equilibrada, com prevalência de projetos em regiões estagnadas (1.162), frente a regiões de alta renda (943). Olhando por setores econômicos, os projetos são em maior parte ligados à indústria de transformação, com cerca de 76\% do total, e à infraestrutura (16\%), distribuição que se repete em nível de superintendência e tipologia.

Em síntese, a estrutura vigente de instrumentos da PNDR combina três mecanismos complementares de intervenção: os Fundos Constitucionais, voltados ao crédito subsidiado de ampla capilaridade territorial e com foco preferencial em micro e pequenos produtores; os Fundos de Desenvolvimento, direcionados a projetos estruturantes de grande porte; e os Incentivos Fiscais, que operam via redução da carga tributária para atrair investimentos privados às regiões prioritárias. Embora compartilhem o objetivo de promover o desenvolvimento das regiões menos dinâmicas, os instrumentos diferem substancialmente em volume de recursos, público-alvo, mecanismo de operação e grau de aderência aos objetivos da política \cite{PortugalSilva2020, Portugal2024}. Essa diversidade institucional, somada à persistência das desigualdades regionais que motivam a intervenção estatal, torna imprescindível a avaliação rigorosa dos efeitos de cada instrumento sobre os indicadores socioeconômicos locais, tarefa que a literatura empírica tem buscado enfrentar com crescente sofisticação metodológica. Para mapear sistematicamente essa produção acadêmica e sintetizar suas contribuições, este artigo emprega uma revisão sistemática da literatura, cujos procedimentos metodológicos são detalhados a seguir.
